\begin{afterword}
\summaryhead

The essay argues that a conclusion is only as good as the process that chose it. A ``clever arguer'' is paid to argue that box B holds a diamond. He writes the conclusion first, at the bottom of the page, and then lists above it only the clues that favor box B.

The author's point is that once the conclusion is written, the arguments added above it cannot change what it is worth. It tells you who paid, not what is in the box. Compare a curious observer who lists the clues for both boxes and only then estimates a probability. That number depends on the evidence.

The same happens inside one head. If you decide your squealing brakes don't need fixing because repairs are expensive, and then look for reasons, the reasons are decoration. The essay ends with a warning: this is a check on your own thinking, not a weapon against other people's conclusions.

\responsehead

The essay makes one familiar and correct point. When someone decides first and argues second, the arguments tell you little. It makes the point with a strong image. Almost everything else in it works against clarity.

The vocabulary persuades where it should explain. Everett branches, Tegmark duplicates, entanglement, algorithms: none of them is needed, because the argument uses nothing beyond ordinary probability over repeated auctions. The words borrow the authority of physics and computing, and they make an old observation about motivated reasoning feel like a technical result. A careful reader should be suspicious when a text brings in heavy machinery to carry a light load.

The central claim, as written, is false. The essay says the arguer's page is evidence ``only of which owner paid the higher bid.'' But the clues on the page are real. A blue stamp is a blue stamp, whoever points it out. What the page hides is the clues left out, so the right response is to discount it, not to discard it. The author corrected this in a post the next day. Yet this is the version LessWrong still includes in its Highlights from the Sequences, and a reader who stops here learns the wrong rule.

The advice also cannot be used. The essay says your success depends on ``whichever algorithm actually writes the bottom line of your thoughts,'' but it gives no way to find out which one that is. From the inside, motivated reasoning feels like honest reasoning; that is the whole problem. A warning you cannot apply to yourself is easiest to apply to others. The last paragraph warns against doing so, but gives no method for checking oneself either.

In short: a good image and a sound instinct, dressed in borrowed jargon, with its central claim overstated and no method for acting on it.
\end{afterword}
